%%%%%%%%%%%%%%%%%%%%%%%%%%%%%%%%%%%%%%%%%%%%%%%%%%%%%%%%%%%%%%%%%%%%%%%%%%%%%%%
% CHAPTER 13: BEHAVIORAL CHANGE JOURNEY (BCJ)
%%%%%%%%%%%%%%%%%%%%%%%%%%%%%%%%%%%%%%%%%%%%%%%%%%%%%%%%%%%%%%%%%%%%%%%%%%%%%%%
% Version: 2.0 (January 2026)
%
% Purpose: Introduces the BCJ as the 8th CORE dimension, answering Question 5:
%          "Where in the Behavioral Change Journey does the decision-maker
%          find themselves at the moment of the decision/behavior?"
%
% Primary Appendix: AW - CORE-STAGE (Behavioral Change Journey)
% Secondary Appendices: AU (CORE-AWARE), AV (CORE-READY), V (CORE-WHEN)
%
% Prerequisites: Ch. 11 (Awareness), Ch. 12 (Willingness)
% Leads to: Ch. 14 (Behavioral Change Segments), Ch. 15 (WEC Synthesis)
%
% Chapter Type: A (CORE Chapter - answers STAGE question)
%%%%%%%%%%%%%%%%%%%%%%%%%%%%%%%%%%%%%%%%%%%%%%%%%%%%%%%%%%%%%%%%%%%%%%%%%%%%%%%

\chapter{Behavioral Change Journey: Stage-Dependent Dynamics}
\label{ch:bcj}

%%%%%%%%%%%%%%%%%%%%%%%%%%%%%%%%%%%%%%%%%%%%%%%%%%%%%%%%%%%%%%%%%%%%%%%%%%%%%%%
% CHAPTER SCOPE BOX
%%%%%%%%%%%%%%%%%%%%%%%%%%%%%%%%%%%%%%%%%%%%%%%%%%%%%%%%%%%%%%%%%%%%%%%%%%%%%%%
\begin{tcolorbox}[
    colback=purple!5!white,
    colframe=purple!75!black,
    title={\textbf{Chapter Scope: Ziel / In-Scope / Out-of-Scope}},
    fonttitle=\bfseries
]

\textbf{Ziel (Objective):}
Establish the formal foundation for stage dynamics, phi phases within the 10C CORE framework.

\vspace{0.3cm}
\textbf{In-Scope:}
\begin{itemize}[nosep]
    \item Core definitions and axioms for stage dynamics
    \item Integration with other 10C dimensions
    \item Formal mathematical foundations
    \item Key theorems and their implications
\end{itemize}

\vspace{0.3cm}
\textbf{Out-of-Scope (delegiert an Appendices):}
\begin{itemize}[nosep]
    \item Detailed proofs and derivations $\rightarrow$ CORE Appendix
    \item Empirical calibration $\rightarrow$ Appendix BBB
    \item Domain-specific applications $\rightarrow$ Application chapters
\end{itemize}

\end{tcolorbox}



%%%%%%%%%%%%%%%%%%%%%%%%%%%%%%%%%%%%%%%%%%%%%%%%%%%%%%%%%%%%%%%%%%%%%%%%%%%%%%%
% QUICK REFERENCE BOX
%%%%%%%%%%%%%%%%%%%%%%%%%%%%%%%%%%%%%%%%%%%%%%%%%%%%%%%%%%%%%%%%%%%%%%%%%%%%%%%
\begin{tcolorbox}[
    colback=gray!5!white,
    colframe=gray!75!black,
    title={\textbf{Begriffe in diesem Kapitel}},
    fonttitle=\bfseries
]
\small
\begin{tabular}{@{}ll@{}}
\textbf{Term} & \textbf{Definition} \\
\midrule
BCJ & Behavioral Change Journey (6 stages) \\
$S_1$--$S_6$ & PreAware → Aware → Triggered → Acting → Maintaining → Transformed \\
Stage Transition & Movement between BCJ stages \\
Boundary Conditions & 152 theorems defining stage boundaries \\
\end{tabular}
\end{tcolorbox}

%==============================================================================
% DIE 9 FRAGEN NAVIGATION BOX
%==============================================================================
\BCMQuestionNav{13}

%==============================================================================
% CORE-STAGE CONNECTION BOX
%==============================================================================
\begin{tcolorbox}[
    colback=teal!5!white,
    colframe=teal!75!black,
    title={\textbf{10C CORE Connection: STAGE (Question 5)}},
    fonttitle=\bfseries
]
\textbf{The 8th CORE Question:} Where in the Behavioral Change Journey?

\medskip
\begin{tabular}{ll}
\textbf{CORE Appendix:} & AW (CORE-STAGE) \\
\textbf{Output:} & Stage $S \in \{S_1, S_2, S_3, S_4, S_5, S_6\}$ \\
\textbf{Key Insight:} & Stage determines which interventions work \\
\textbf{Novel Contribution:} & 152 boundary conditions for stage-behavior mapping
\end{tabular}

\medskip
\textit{This chapter extends 7C to 10C by adding STAGE as the temporal dimension of behavioral change.}
\end{tcolorbox}

%==============================================================================
% APPENDIX REFERENCES BOX
%==============================================================================
\begin{tcolorbox}[
    colback=yellow!10!white,
    colframe=orange!75!black,
    title={\textbf{Appendix References for This Chapter}},
    fonttitle=\bfseries
]
\begin{tabular}{lll}
\textbf{Appendix} & \textbf{Content} & \textbf{Section} \\
\midrule
AW (CORE-STAGE) & 85 Axioms, 152 Theorems & All sections \\
AU (CORE-AWARE) & Awareness function $A(\cdot)$ & \ref{sec:bcj-awareness} \\
AV (CORE-READY) & Willingness threshold $\theta$ & \ref{sec:bcj-willingness} \\
V (CORE-WHEN) & Context modulation $\Psi$ & \ref{sec:bcj-context} \\
F (REF-EXAMPLES) & Worked examples & \ref{sec:bcj-applications}
\end{tabular}
\end{tcolorbox}

%==============================================================================
% CROSS-REFERENCE: EIT-EMP / EMERGE ALGORITHM
%==============================================================================
\begin{tcolorbox}[
    colback=blue!5!white,
    colframe=blue!75!black,
    title={\textbf{Cross-Reference: EMERGE Algorithm (Chapter 11, EIT-EMP)}},
    fonttitle=\bfseries
]
\textbf{STAGE $\rightarrow$ Intervention Design:} Die BCJ-Stage $S \in \{S_1, ..., S_6\}$ bestimmt den Stage-Faktor $f_S$ im EMERGE-Algorithmus (Chapter 11) und damit die optimale Portfoliogröße $K^*$.

\medskip
\textbf{Stage $\rightarrow$ $f_S$ Mapping:}
\begin{center}
\begin{tabular}{llcl}
\textbf{BCJ Stage} & \textbf{Charakteristik} & \textbf{$f_S$} & \textbf{Intervention-Fokus} \\
\midrule
$S_1$ (Unaware) & $AWX \approx 0$ & 0.6 & Awareness schaffen \\
$S_2$ (Aware) & $AWX > 0$, abstrakt & 0.7 & Persönliche Relevanz \\
$S_3$ (Considering) & $WAX < \theta$ & 1.0 & Barrieren adressieren \\
$S_4$ (Intending) & $WAX \approx \theta$ & 1.2 & Intention-Behavior Gap schließen \\
$S_5$ (Acting) & $WAX > \theta$ & 1.0 & Friktion minimieren \\
$S_6$ (Maintaining) & $\theta \to 0$ & 0.8 & Habit-Lock-In sichern \\
\end{tabular}
\end{center}

\medskip
\textbf{EMERGE-Formel (EIT-EMP-2, EXC-3 Hybrid):}
\[
K^* = K_{\text{base}} \cdot f_B + \delta_C + \boxed{\delta_S} + \delta_\gamma + \delta_L
\]
Nur $f_B$ ist multiplikativ. Der Stage-Adjustment $\delta_S$ ist additiv und maximal in $S_4$ (Intending), wo der kritische Übergang zur Handlung stattfindet.

\medskip
\textbf{Verbindung zu Phase-Affinity (Chapter 18):}
Die Phase-Affinity Matrix $\alpha_{ij}$ (Chapter 18) quantifiziert, welche Interventionstypen in welcher BCJ-Stage wirksam sind. $f_S$ aggregiert diese Information für die $K^*$-Berechnung.

\medskip
\textbf{Workflow:} Chapter 13 (STAGE: $S$) $\rightarrow$ Chapter 11 (EMERGE: $f_S$, $K^*$) $\rightarrow$ Chapters 17--22 (Implementation)

\medskip
\textit{Siehe Chapter 11, Section 11.5 für die vollständige EMERGE-Spezifikation.}
\end{tcolorbox}


%%%%%%%%%%%%%%%%%%%%%%%%%%%%%%%%%%%%%%%%%%%%%%%%%%%%%%%%%%%%%%%%%%%%%%%%%%%%%%%
%                          PART I: THE JOURNEY PERSPECTIVE
%%%%%%%%%%%%%%%%%%%%%%%%%%%%%%%%%%%%%%%%%%%%%%%%%%%%%%%%%%%%%%%%%%%%%%%%%%%%%%%

\section{Introduction: From Static to Dynamic}
\label{sec:bcj-intro}

%------------------------------------------------------------------------------

Behavioral change is not a single event but a \textit{process} that unfolds over time.
The \textbf{Behavioral Change Journey (BCJ)} captures where an individual stands in
this process at any given moment.

\begin{intuition}[Why Stages Matter]
Imagine two people both aware that exercise is healthy:
\begin{itemize}
    \item \textbf{Maria:} "I know I should exercise, but I haven't really thought about it."
    \item \textbf{Thomas:} "I've planned to start running next Monday. I even bought shoes."
\end{itemize}
Both have awareness ($AWX > 0$), but they are at \textit{different stages}.
An intervention that works for Thomas (e.g., a running buddy) may fail for Maria
(who first needs to see personal relevance). \textbf{Stage determines intervention effectiveness.}
\end{intuition}

\paragraph{The Central Question.}
This chapter addresses the 5th of the 9 BCM Questions:

\begin{center}
\fcolorbox{teal!75!black}{teal!5!white}{\parbox{0.85\textwidth}{
\textbf{Question 5:} Where in the \textbf{Behavioral Change Journey}
does the decision-maker find themselves at the moment of the decision/behavior?
}}
\end{center}

\paragraph{Chapter Overview.}
This chapter proceeds as follows:
\begin{itemize}
    \item Section \ref{sec:bcj-stages}: Defines the six stages of the BCJ
    \item Section \ref{sec:bcj-formal}: Formalizes BCJ as a modulator of AWX and $\theta$
    \item Section \ref{sec:bcj-transitions}: Specifies transition dynamics
    \item Section \ref{sec:bcj-integration}: Connects BCJ to the 10C CORE framework
    \item Section \ref{sec:bcj-applications}: Provides worked examples
    \item Section \ref{sec:bcj-summary}: Summarizes key contributions
\end{itemize}


%%%%%%%%%%%%%%%%%%%%%%%%%%%%%%%%%%%%%%%%%%%%%%%%%%%%%%%%%%%%%%%%%%%%%%%%%%%%%%%
%                          PART II: THE SIX STAGES
%%%%%%%%%%%%%%%%%%%%%%%%%%%%%%%%%%%%%%%%%%%%%%%%%%%%%%%%%%%%%%%%%%%%%%%%%%%%%%%

\section{The Six Stages of the BCJ}
\label{sec:bcj-stages}

%------------------------------------------------------------------------------

\begin{definition}[Behavioral Change Journey]
The Behavioral Change Journey (BCJ) is an ordered sequence of six stages
$S = \{S_1, S_2, S_3, S_4, S_5, S_6\}$ that characterize an individual's position
in the process of behavioral change. Each stage is defined by characteristic
values of Awareness ($AWX$), Willingness ($WAX$), and Threshold ($\theta$).
\end{definition}

\subsection{Stage 1: Unaware ($S_1$)}
\label{sec:bcj-s1}

\begin{definition}[Stage 1: Unaware]
No awareness of the behavior or its personal relevance.
\begin{equation}
S_1: \quad AWX \approx 0, \quad \theta \to \infty
\end{equation}
\end{definition}

\paragraph{Characteristics.}
\begin{itemize}
    \item The behavior option is not in the consideration set
    \item No mental representation of costs/benefits
    \item Threshold effectively infinite (behavior cannot be triggered)
\end{itemize}

\paragraph{Typical Statement.}
\textit{"Heat pump? Never heard of it. What's that?"}

\paragraph{Intervention Focus.}
Creating basic awareness through exposure, information, salience.

\begin{cross-reference}
See Appendix AWA (CORE-AWARE) for the formal Awareness function $A(\cdot)$
that determines transition from $S_1$ to $S_2$.
\end{cross-reference}

%------------------------------------------------------------------------------
\subsection{Stage 2: Aware ($S_2$)}
\label{sec:bcj-s2}

\begin{definition}[Stage 2: Aware]
Knows the option exists but does not perceive personal relevance.
\begin{equation}
S_2: \quad AWX > 0 \text{ (low)}, \quad \theta \text{ high}
\end{equation}
\end{definition}

\paragraph{Characteristics.}
\begin{itemize}
    \item The option enters the mental consideration set
    \item Abstract knowledge without personal connection
    \item "That's a thing other people do"
\end{itemize}

\paragraph{Typical Statement.}
\textit{"Yes, heat pumps exist. I guess they're somehow relevant for the environment."}

\paragraph{Intervention Focus.}
Establishing personal relevance: "This applies to YOU."

%------------------------------------------------------------------------------
\subsection{Stage 3: Considering ($S_3$)}
\label{sec:bcj-s3}

\begin{definition}[Stage 3: Considering]
Actively evaluating the option for personal adoption.
\begin{equation}
S_3: \quad AWX \text{ medium}, \quad WAX < \theta
\end{equation}
\end{definition}

\paragraph{Characteristics.}
\begin{itemize}
    \item Active information search
    \item Cost-benefit analysis in progress
    \item Willingness growing but not yet sufficient
\end{itemize}

\paragraph{Typical Statement.}
\textit{"Hmm, maybe I should look into getting a heat pump..."}

\paragraph{Intervention Focus.}
Concretizing benefits, reducing uncertainty, addressing specific concerns.

%------------------------------------------------------------------------------
\subsection{Stage 4: Intending ($S_4$)}
\label{sec:bcj-s4}

\begin{definition}[Stage 4: Intending]
Has formed an intention to act but has not yet acted.
\begin{equation}
S_4: \quad WAX \approx \theta, \quad \text{Intention-Behavior Gap}
\end{equation}
\end{definition}

\paragraph{Characteristics.}
\begin{itemize}
    \item Clear intention exists
    \item "I want to do this" but "not yet"
    \item Threshold and willingness approximately equal
\end{itemize}

\paragraph{Typical Statement.}
\textit{"Yes, I definitely want a heat pump. I'll get to it... someday."}

\paragraph{Intervention Focus.}
Lowering $\theta$ through: commitment devices, implementation intentions, defaults.

\begin{cross-reference}
See Appendix REA (CORE-READY) for the formal specification of $\theta$
and the conditions under which $WAX > \theta$ triggers action.
\end{cross-reference}

%------------------------------------------------------------------------------
\subsection{Stage 5: Acting ($S_5$)}
\label{sec:bcj-s5}

\begin{definition}[Stage 5: Acting]
Actively performing the target behavior.
\begin{equation}
S_5: \quad WAX > \theta
\end{equation}
\end{definition}

\paragraph{Characteristics.}
\begin{itemize}
    \item Behavior is being executed
    \item Decision made, action in progress
    \item Still requires conscious effort
\end{itemize}

\paragraph{Typical Statement.}
\textit{"I'm installing the heat pump right now!"}

\paragraph{Intervention Focus.}
Minimizing friction, simplifying processes, reducing implementation barriers.

%------------------------------------------------------------------------------
\subsection{Stage 6: Maintaining / Habit ($S_6$)}
\label{sec:bcj-s6}

\begin{definition}[Stage 6: Maintaining]
The behavior has become automatic (habituated).
\begin{equation}
S_6: \quad \theta \to 0, \quad \text{Automaticity achieved}
\end{equation}
\end{definition}

\paragraph{Characteristics.}
\begin{itemize}
    \item No conscious decision required
    \item Behavior runs on autopilot
    \item New identity: "I am someone who does this"
\end{itemize}

\paragraph{Typical Statement.}
\textit{"Heat pump? Of course, it's just what we have. I don't think about it."}

\paragraph{Intervention Focus.}
Feedback, reinforcement, preventing relapse.

\begin{intuition}[The Habit Lock-In]
Once $\theta \to 0$, the behavior becomes "locked in." This is the
\textbf{Habit Lock-In Time} ($T_{lock}$), a novel prediction of EBF:
the behavior continues with minimal psychological cost, creating behavioral
inertia that resists change.
\end{intuition}


%%%%%%%%%%%%%%%%%%%%%%%%%%%%%%%%%%%%%%%%%%%%%%%%%%%%%%%%%%%%%%%%%%%%%%%%%%%%%%%
%                     PART III: FORMAL SPECIFICATION
%%%%%%%%%%%%%%%%%%%%%%%%%%%%%%%%%%%%%%%%%%%%%%%%%%%%%%%%%%%%%%%%%%%%%%%%%%%%%%%

\section{Formal Specification: BCJ as Modulator}
\label{sec:bcj-formal}

%------------------------------------------------------------------------------

\subsection{Stage Index Definition}

\begin{definition}[BCJ Stage Index]
Let $S \in \{1, 2, 3, 4, 5, 6\}$ denote the current BCJ stage:
\begin{equation}
S = \begin{cases}
1 & \text{Unaware} \\
2 & \text{Aware} \\
3 & \text{Considering} \\
4 & \text{Intending} \\
5 & \text{Acting} \\
6 & \text{Maintaining}
\end{cases}
\label{eq:bcj-stage-index}
\end{equation}
\end{definition}

\subsection{Stage-Dependent Parameter Modulation}

Each stage implies characteristic ranges for the core BCM parameters:

\begin{table}[h]
\centering
\begin{tabular}{lcccl}
\toprule
\textbf{Stage} & $\varphi$ & $AWX$ & $\theta$ & \textbf{Primary Focus} \\
\midrule
$S_1$ (Unaware) & -- & $\approx 0$ & $\to \infty$ & Create awareness \\
$S_2$ (Aware) & low & low & high & Establish relevance \\
$S_3$ (Considering) & variable & medium & medium & Concretize benefits \\
$S_4$ (Intending) & rising & high & medium & Lower threshold \\
$S_5$ (Acting) & high & high & low & Eliminate friction \\
$S_6$ (Maintaining) & -- & -- & $\to 0$ & Reinforce habit \\
\bottomrule
\end{tabular}
\caption{BCJ Stage characteristics and intervention focus}
\label{tab:bcj-stages}
\end{table}

\subsection{The Stage Modulation Function}

\begin{definition}[Stage Modulation]
The effective threshold $\theta^{eff}$ is modulated by stage $S$:
\begin{equation}
\theta^{eff}(S) = \theta_0 \cdot m(S)
\label{eq:theta-modulation}
\end{equation}
where $m(S)$ is the stage modulation factor with $m(S_6) \to 0$ and $m(S_1) \to \infty$.
\end{definition}


%%%%%%%%%%%%%%%%%%%%%%%%%%%%%%%%%%%%%%%%%%%%%%%%%%%%%%%%%%%%%%%%%%%%%%%%%%%%%%%
%                      PART IV: TRANSITION DYNAMICS
%%%%%%%%%%%%%%%%%%%%%%%%%%%%%%%%%%%%%%%%%%%%%%%%%%%%%%%%%%%%%%%%%%%%%%%%%%%%%%%

\section{Transition Dynamics}
\label{sec:bcj-transitions}

%------------------------------------------------------------------------------

\subsection{Forward Transitions}

\begin{definition}[Stage Transition Probability]
The probability of transitioning from stage $S_i$ to $S_{i+1}$ depends on:
\begin{equation}
P(S_i \to S_{i+1}) = f(AWX, WAX, \theta, \Psi, \text{Intervention})
\label{eq:transition-prob}
\end{equation}
where $\Psi$ is the context vector.
\end{definition}

\paragraph{Transition Conditions.}
\begin{itemize}
    \item $S_1 \to S_2$: Requires $AWX$ crossing awareness threshold
    \item $S_2 \to S_3$: Requires personal relevance recognition
    \item $S_3 \to S_4$: Requires intention formation ($WAX$ approaching $\theta$)
    \item $S_4 \to S_5$: Requires $WAX > \theta$ (action trigger)
    \item $S_5 \to S_6$: Requires repeated action until automaticity
\end{itemize}

\subsection{Backward Transitions (Relapse)}
\label{sec:bcj-relapse}

\begin{intuition}[Relapse as Stage Regression]
Behavioral change is not unidirectional. Individuals can regress to earlier stages:
\begin{itemize}
    \item $S_6 \to S_5$: Habit disruption (new context breaks automaticity)
    \item $S_5 \to S_4$: Action failure (implementation barriers)
    \item $S_4 \to S_3$: Intention decay (motivation fades)
\end{itemize}
\end{intuition}

\begin{definition}[Relapse Probability]
\begin{equation}
P(S_i \to S_{i-1}) = g(\Delta\Psi, \text{Time since action}, \text{Alternative options})
\label{eq:relapse-prob}
\end{equation}
\end{definition}

\begin{cross-reference}
Appendix STA (CORE-STAGE) contains 152 boundary conditions specifying
exact conditions for forward and backward transitions across 29 model categories.
\end{cross-reference}


%%%%%%%%%%%%%%%%%%%%%%%%%%%%%%%%%%%%%%%%%%%%%%%%%%%%%%%%%%%%%%%%%%%%%%%%%%%%%%%
%                     PART V: INTEGRATION WITH 7C/10C
%%%%%%%%%%%%%%%%%%%%%%%%%%%%%%%%%%%%%%%%%%%%%%%%%%%%%%%%%%%%%%%%%%%%%%%%%%%%%%%

\section{Integration with the CORE Framework}
\label{sec:bcj-integration}

%------------------------------------------------------------------------------

The BCJ extends the 7C framework to 10C by adding \textbf{STAGE} as the
temporal dimension of behavioral change.

\subsection{Connection to CORE-AWARE (Chapter 11)}
\label{sec:bcj-awareness}

\begin{cross-reference}
The Awareness function $A(\cdot)$ from CORE-AWARE determines:
\begin{itemize}
    \item Initial entry into $S_2$ (becoming aware)
    \item Quality of awareness that enables $S_2 \to S_3$
    \item Awareness decay that can cause regression
\end{itemize}
See Chapter 11 and Appendix AWA for full specification.
\end{cross-reference}

\subsection{Connection to CORE-READY (Chapter 12)}
\label{sec:bcj-willingness}

\begin{cross-reference}
The Willingness function $WAX$ and threshold $\theta$ from CORE-READY determine:
\begin{itemize}
    \item The $WAX \approx \theta$ condition defining $S_4$
    \item The $WAX > \theta$ trigger for $S_5$
    \item The $\theta \to 0$ habituation in $S_6$
\end{itemize}
See Chapter 12 and Appendix REA for full specification.
\end{cross-reference}

\subsection{Connection to CORE-WHEN (Context)}
\label{sec:bcj-context}

\begin{cross-reference}
The context vector $\Psi$ from CORE-WHEN modulates all BCJ parameters:
\begin{itemize}
    \item $\Psi$ affects baseline $\theta_0$ at each stage
    \item Context changes can trigger stage transitions
    \item $\Delta\Psi$ is a primary driver of relapse
\end{itemize}
See Appendix WEN (CORE-WHEN) for the 8-dimensional context specification.
\end{cross-reference}

\subsection{The 10C Integration Table}

\begin{table}[h]
\centering
\begin{tabular}{lll}
\toprule
\textbf{CORE} & \textbf{Question} & \textbf{BCJ Role} \\
\midrule
WHO (AAA) & Who has utility? & Stage affects all levels $L$ \\
WHAT (C) & What is utility? & Stage modulates FEPSDE weights \\
HOW (B) & How interact? & Stage affects complementarity $\gamma$ \\
WHEN (V) & When context? & Context modulates stage transitions \\
WHERE (BBB) & Where parameters? & Stage determines parameter ranges \\
AWARE (AU) & How aware? & Awareness enables early stages \\
READY (AV) & How willing? & Willingness determines action stages \\
\textbf{STAGE (AW)} & \textbf{Where in journey?} & \textbf{Stage determines all} \\
\bottomrule
\end{tabular}
\caption{Integration of STAGE (8th CORE) with the complete 10C framework}
\label{tab:10c-integration}
\end{table}


%%%%%%%%%%%%%%%%%%%%%%%%%%%%%%%%%%%%%%%%%%%%%%%%%%%%%%%%%%%%%%%%%%%%%%%%%%%%%%%
%                     PART VI: WORKED EXAMPLES
%%%%%%%%%%%%%%%%%%%%%%%%%%%%%%%%%%%%%%%%%%%%%%%%%%%%%%%%%%%%%%%%%%%%%%%%%%%%%%%

\section{Worked Examples}
\label{sec:bcj-applications}

%------------------------------------------------------------------------------

\subsection{Example 1: Anna's Heat Pump Journey}

\paragraph{Baseline: Stage $S_3$ (Considering).}
\begin{itemize}
    \item Statement: "Hmm, heat pump... should probably look into it..."
    \item Parameters: $AWX = 0.43$, $WAX = 0.48$, $\theta = 1.32$
    \item Analysis: $WAX < \theta$ $\Rightarrow$ not ready to act
\end{itemize}

\paragraph{Intervention: Social Proof + Subsidy Information.}
\begin{itemize}
    \item "The Müllers next door installed one. 40\% subsidy available!"
    \item Effect: $AWX \uparrow$, $\theta \downarrow$ (relevance + barrier reduction)
\end{itemize}

\paragraph{After Intervention: Stage $S_5$ (Acting).}
\begin{itemize}
    \item Statement: "Yes, I want this! Let's get a quote!"
    \item Parameters: $AWX = 0.87$, $WAX = 0.90$, $\theta = 0.76$
    \item Analysis: $WAX > \theta$ $\Rightarrow$ action triggered
\end{itemize}

\begin{equation}
\Delta S = S_5 - S_3 = 2 \text{ stages forward (skipped } S_4\text{)}
\label{eq:anna-transition}
\end{equation}

\subsection{Example 2: Multi-Level BCJ (Organizational)}

\begin{intuition}[BCJ at Different Levels]
The BCJ applies not only to individuals but to all levels $L$:
\begin{itemize}
    \item $L=0$ (Individual): Anna considering a heat pump
    \item $L=1$ (Household): Family deciding on solar panels
    \item $L=2$ (Organization): Company adopting sustainability policy
    \item $L=\infty$ (Society): Nation transitioning to renewable energy
\end{itemize}
Each level has its own BCJ with level-specific transition dynamics.
\end{intuition}

\begin{cross-reference}
Appendix STA contains boundary conditions for multi-level BCJ from
models across 29 categories including Matching Markets (T133-T137),
Information Economics (T138-T142), and Development Economics (T148-T152).
\end{cross-reference}


%%%%%%%%%%%%%%%%%%%%%%%%%%%%%%%%%%%%%%%%%%%%%%%%%%%%%%%%%%%%%%%%%%%%%%%%%%%%%%%
%                          SUMMARY AND TRANSITION
%%%%%%%%%%%%%%%%%%%%%%%%%%%%%%%%%%%%%%%%%%%%%%%%%%%%%%%%%%%%%%%%%%%%%%%%%%%%%%%

\section{Summary}
\label{sec:bcj-summary}

%------------------------------------------------------------------------------

This chapter has established:

\begin{enumerate}
    \item \textbf{BCJ as 8th CORE:} Stage (STAGE) extends 7C to 10C, answering
          "Where in the journey?" as a fundamental dimension of behavioral change.

    \item \textbf{Six Stages:} $S_1$ (Unaware) through $S_6$ (Maintaining), each
          with characteristic AWX, WAX, and $\theta$ profiles.

    \item \textbf{Stage-Dependent Intervention:} Different interventions are
          effective at different stages; one-size-fits-all fails.

    \item \textbf{Transition Dynamics:} Forward transitions require specific
          conditions; backward transitions (relapse) are possible.

    \item \textbf{10C Integration:} BCJ connects to all other CORE dimensions,
          modulating their effects through the stage lens.
\end{enumerate}

%------------------------------------------------------------------------------
\paragraph{Formal Details.}
The complete formal treatment appears in Appendix~AW (CORE-STAGE), which contains:
\begin{itemize}
    \item 85 axioms specifying the BCJ structure
    \item 152 boundary conditions from 29 model categories
    \item Proof of endogenous stage emergence
    \item The $T_{lock}$ (Habit Lock-In Time) prediction
\end{itemize}

%------------------------------------------------------------------------------
\paragraph{What Comes Next.}
Chapter 14 introduces \textbf{Behavioral Change Segments (BCS)}, which combine
BCJ stages with individual characteristics to create actionable population
segments for targeted intervention design.

Then Chapter 15 integrates BCJ (journey stage) with BCS (segment type) to answer:
**HOW do we structure coherent multi-layer interventions?** Using the HI CORE-HIERARCHY
framework, Chapter 15 shows that decision hierarchy complexity depends on complementarity
and modularity---captured by the universal formula $N_{L2} = \alpha \cdot \gamma_{avg} \times n \times (1-m) / \log(n)$.

\begin{cross-reference}
\textbf{Reading Path:}
\begin{itemize}
    \item For formal BCJ specification: Appendix STA (CORE-STAGE)
    \item For awareness mechanics: Chapter 11 + Appendix AWA
    \item For willingness mechanics: Chapter 12 + Appendix REA
    \item For segmentation: Chapter 14 (Behavioral Change Segments)
    \item For coherent intervention design: Chapter 15 (Decision Hierarchy via HI CORE-HIERARCHY)
    \item For probability prediction: Chapter 16 (Effective Probability of Change)
    \item \textbf{For intervention optimization: Chapter 11 (EMERGE Algorithm) + Appendix IE (EIT-EMP)}
\end{itemize}
\end{cross-reference}

%------------------------------------------------------------------------------
\paragraph{Connection to Intervention Design (EIT-EMP).}
The BCJ stage $S$ is a primary input for the EMERGE algorithm (Chapter 11). Each stage
maps to a stage factor $f_S$ that modulates the optimal portfolio size $K^*$: early stages
($S_1$, $S_2$) have lower $f_S$ (fewer interventions needed for awareness), the critical
$S_4$ (Intending) has maximum $f_S = 1.2$ (more interventions to close the intention-behavior gap),
and $S_6$ (Maintaining) has reduced $f_S = 0.8$ (habit lock-in requires less active intervention).
The Phase-Affinity Matrix $\alpha_{ij}$ (Chapter 18) provides the detailed mapping of which
intervention types are effective at each BCJ stage.

%%%%%%%%%%%%%%%%%%%%%%%%%%%%%%%%%%%%%%%%%%%%%%%%%%%%%%%%%%%%%%%%%%%%%%%%%%%%%%%
% END OF CHAPTER 13
%%%%%%%%%%%%%%%%%%%%%%%%%%%%%%%%%%%%%%%%%%%%%%%%%%%%%%%%%%%%%%%%%%%%%%%%%%%%%%%
